\chapter*{读书报告}
\addcontentsline{toc}{chapter}{读书报告}
\markboth{读书报告}{读书报告}
\begin{center}
\large 从辛几何到可计算的 $A_\infty$ 范畴
\end{center}

\begin{quote}
\small 本报告基于全书通读与本次校订工作写成。文中页码均为原书印刷页码；定理编号从原书。书末译者注共 187 条，可与本报告互参。
\end{quote}

\section*{一、这本书要解决什么问题}

Fukaya 范畴把辛流形的几何编码成一个 $A_\infty$ 范畴：对象是 Lagrangian 子流形，态射是 Floer 上链群，复合映射由伪全纯圆盘的计数给出。这个定义在概念上很自然，在计算上却几乎无法直接使用——对象本身就无法分类，更不用说数清楚所有圆盘。Seidel 这本书的目标，不是再讲一遍定义，而是回答一个更实际的问题：\emph{在什么条件下，Fukaya 范畴可以被约化为有限、可计算的数据？}

他的答案是：把辛流形限制到正合（exact）且 Calabi--Yau 的仿射代数簇，用 Picard--Lefschetz 理论把几何分解为消失圈与辫单值变换，再用 $A_\infty$ 范畴的代数形式（扭曲复形、分裂闭包、变异）把归纳计算组织起来。书中的定理 A、B 与推论 C（原书 pp.~4--5）是这条路线的主定理；全书三章分别准备代数、分析、几何三套工具。

\section*{二、三章的结构}

\textbf{第一章（§1--6）是纯同调代数。}它从 $A_\infty$ 范畴的定义讲起，重点不在铺陈理论，而在几件后文真正用得到的事：恒等态射在 $A_\infty$ 层次并不自动存在；映射锥与正合三角如何构造；幂等元如何形式地分裂出直和项；以及“沿一个对象的扭转”如何实现变异。作者明确说，这一章只详细讨论恒等态射，因为它在 Fukaya 范畴中恰恰最棘手。章末引入 $\Z/2$ 等变结构，为第三章的二重覆盖做准备。

\textbf{第二章（§7--14）是 Floer 理论。}它先以一点辛几何（§7）作玩具模型，再回顾经典 Floer 上同调（§8），然后逐步给出正合 Fukaya 范畴的“初步版本”（§9）与“完整版本”（§12）。中间三节分别处理基本性质（§10）、指标与行列式线（§11）、曲面上的多边形（§13）。分次（grading）与 Pin 结构是这一章的技术核心：前者把 Floer 群变成 $\Z$ 分次对象，后者决定符号，因而只在 $\operatorname{char}\K\ne2$ 时必需。§14 讨论带辛对合的流形及其等变 Fukaya 范畴，是第三章定义 $\F(\pi)$ 的直接工具。

\textbf{第三章（§15--20）是 Picard--Lefschetz 理论。}它从 Lefschetz 纤维化的初等几何开始（§15），引入消失圈、Lefschetz 顶针与匹配圈（§16），用伪全纯截面证明扭转函子与辛 Dehn 扭转的相容性（§17），再借二重覆盖定义 Lefschetz 纤维化的 Fukaya 范畴 $\F(\pi)$（§18）。§19 回到仿射簇，证明引言中的三个定理；§20 用 $(A_m)$ 型 Milnor 纤维作例子，展示这套机器如何真正算出东西。

\section*{三、第一章：把“算”变成“形式化”}

$A_\infty$ 范畴的定义是一组多重线性映射 $\mu^d$，满足带符号的二次关系。与 dg 范畴不同，它不要求 $\mu^1$ 与 $\mu^2$ 先满足结合律；结合律只在同调层次成立。这带来两个后果。其一，恒等态射通常不是严格存在的，只能在一个拟同构的模型里找到；Seidel 用同伦扰动（homological perturbation）处理这一点，并反复提醒读者：在 $A_\infty$ 层次做计算时，不能随意把恒等态射当成真的单位元。其二，映射锥不能按朴素方式定义，必须引入扭曲复形。

扭曲复形的对象是一族对象 $X_i$ 加上满足 Maurer--Cartan 方程的矩阵 $\delta=(\delta_{ij})$，$\delta_{ij}\in\hom^1(X_i,X_j)$。它的上同调 $H^0(\Tw\A)$ 是三角范畴；这一步把“映射锥”“正合三角”从链复形的直觉搬到了 $A_\infty$ 世界。接下来是分裂闭包：把每个幂等自同态的像形式地加进来，得到 $\Pi(\Tw\A)$。作者强调，分裂生成（split-generation）是很弱的要求——它甚至不保证 $K_0$ 由这些生成元的类张成——但正因为它弱，才可能在几何中实现。

扭转（twist）是本章与几何的接口。给定对象 $E$，沿 $E$ 的扭转 $T_E$ 在三角范畴中对应“变异”（mutation）；在 $A_\infty$ 层次，它由一串带符号的复合映射给出。§5 证明 $T_E$ 是 $A_\infty$ 函子，并在 $E$ 为球面对象时给出逆。这条代数线索在第三章会变成：沿消失圈的扭转等于沿匹配圈的 Dehn 扭转。

\section*{四、第二章：正合 Floer 理论的技术骨架}

第二章的出发点是正合辛流形 $(M,\omega=d\theta)$。正合性排除了闭流形，但作者施加凸性（convexity）条件，使有限面积的伪全纯曲线不会跑到“无穷远”。在这个框架里，正合 Lagrangian 子流形 $L$（即 $\theta|_L$ 恰当）没有障碍（unobstructed），每个都给出非零对象，且自同态代数就是经典上同调 $H^*(L)$，没有量子修正。这使 Fukaya 范畴摆脱 Novikov 参数，代价是失去闭流形上的许多例子。

分次是第二个技术支柱。Kontsevich 的分次 Lagrangian 是指一个实值函数 $\alpha_L:L\to\R$，使 $e^{2\pi i\alpha_L}$ 等于某个行列式线的相位；这要求 $2c_1(M)=0$。分次把 Floer 群变成 $\Z$ 分次向量空间，使映射锥、扭转、变异都能逐次进行。Pin 结构则负责符号：伪全纯圆盘模空间的定向由 $L$ 的 Pin 结构决定，若 $\operatorname{char}\K=2$，符号消失，Pin 结构可以省略。

完整定义（§12）用“带边界标记点的圆盘”统一 $\mu^d$：$d$ 个边界标记点把圆周分成 $d$ 段，每段落在一个 Lagrangian 上；模空间是这些圆盘的等价类。这里的技术难点是横截性（transversality）与紧性（compactness）；作者用零阶非齐次项扰动方程，避免虚基本链，但把分析细节压到最低，只给证明概要。§13 讨论曲面：此时伪全纯曲线可以用多边形组合地描述，Floer 理论退化为组合计数；这是全书唯一一处能把分析完全拿掉的地方。

\section*{五、第三章：从几何到辫群，再回到代数}

Picard--Lefschetz 理论的核心图像是：一个 Lefschetz 纤维化 $\pi:E\to\C$ 的临界值绕一圈，纤维上的消失圈就经历一次 Dehn 扭转。把这些扭转按临界值的顺序复合起来，得到辫单值变换（braid monodromy）。Seidel 的观察是：\emph{临界值的 Hurwitz 移动（交换两个临界值的顺序）对应消失圈变异}。这正是定理 17.20（原书 p.~263）的内容，也是“为什么 Fukaya 范畴可以按维数归纳计算”的几何解释；定理 18.24（原书 p.~287）与推论 18.27（原书 p.~291）把它组织成可用的代数形式。

为了把这个观察变成定理，需要两件工具。第一是伪全纯截面：Lefschetz 纤维化的全纯截面在纤维上留下一个点，绕临界值时这个点沿着消失圈移动；计数经过两个截面的伪全纯截面，就得到 Floer 上同调中的扭转映射。第二是 Lefschetz 纤维化的 Fukaya 范畴 $\F(\pi)$：取沿一个光滑纤维分歧的二重覆盖，用 $\Z/2$ 等变构造把纤维与总空间的 Fukaya 范畴联系起来。这个定义不自然，但可以让第二章的等变技术直接复用。

§19 把上述工具用于仿射簇 $X\subset\C^N$ 及其超平面截面 $X^1$。对 $X^1$，消失圈给出自然的一族对象；对 $X$，则用 Lefschetz 束的匹配圈。组合数据（临界值、辫、消失圈的交叠数）编码成一个有限的、有向的 $A_\infty$ 范畴 $\mathcal B$。定理说：$\Pi(\Tw\F(X^1))$ 与 $\Pi(\Tw\mathcal B)$ 拟等价（定理 A）；$\Tw\mathcal B$ 含有与 $\F(X)$ 拟等价的全子范畴（定理 B）；因此 $\F(X)$ 与一个 hom 空间有限维的 dg 范畴拟同构（推论 C）。条件 $d\ne2$、$n\ge1$、$\operatorname{char}\K\ne2$ 都是这条证明路线的技术限制，作者在引言中逐条说明。

这里有一处原书自身的不统一：定理 A 印作 $n\ge1$，但紧接着的解释段（p.~5）与定理 B 之后的说明都按 $n>1$ 使用，并说 $n=1$ 的情形“只是为方便起见”才排除。译文从原书定理陈述，取 $n\ge1$，并在译者注中登记；第 19 节的归纳证明本身也是按 $n\ge1$ 陈述的。

\section*{六、贡献、局限与后续}

这本书的贡献不在于发明 Fukaya 范畴，而在于\emph{把“计算”变成一个有结构的程序}。它给出了从辛几何到组合数据的清晰路线：正合性消去分析上的无穷远，分次与 Pin 结构消去符号歧义，扭曲复形与分裂闭包消去 $A_\infty$ 层次的技术麻烦，PL 理论则把几何归纳为辫群。Hurwitz 移动等于变异这一定理，后来成为镜像对称中“例外集合变异”与“球面对象”理论的基石之一。

局限也很明显。第一，全书只处理正合、Calabi--Yau 的仿射簇；闭流形、一般辛流形不在范围内。第二，分析细节被大幅压缩，读者若想真正补全横截性与紧性，必须回到 Fukaya--Oh--Ohta--Ono 的巨著。第三，条件 $\operatorname{char}\K\ne2$ 与 $d\ne2$ 是证明的技术产物，作者自己说前者“很可能”可以去掉，后者则更本质。第四，除了 $(A_m)$ Milnor 纤维，书中几乎没有例子；计算复杂度在实际中仍然很高。

从今天的视角看，这本书是“Fukaya 范畴的可计算性”这一方向的早期纲领。它之后的 wrapped Fukaya 范畴、family Floer 理论、microlocal sheaf 与 quiver-with-potential 的工作，都可以看作沿着它提出的问题继续走。读这本书的正确方式，不是把它当入门教材逐页啃，而是把它当一张施工图：先看引言与 §19 的定理陈述，再按需回到 §1--6（代数）、§8--12（分析）、§15--18（几何）取工具。

\section*{七、阅读建议}

\begin{itemize}
\item 第一次读：引言（原书 pp.~1--7）$\to$ §19 的定理与证明梗概 $\to$ §20 的例子。先建立“为什么可以归纳计算”的直觉。
\item 第二次读：§1--5（$A_\infty$ 形式化）$\to$ §8--12（Floer 理论）$\to$ §15--18（PL 理论）。每一步都对照书末的符号表与索引。
\item 公式核对：本书的指标与符号约定很密，建议把 §11（指标与行列式线）与 §17（伪全纯截面）的公式单独抄一遍，确认 $(-1)$ 符号与分次移位。
\item 配合阅读：作者在引言中反复引用的 [60]（Fukaya--Oh--Ohta--Ono）是分析细节的正式来源；[131]（Seidel 的 Lefschetz fibration 论文）是第三章的直接前身。
\end{itemize}
